\documentclass[11pt,compress,t,notes=noshow, aspectratio=169, xcolor=table]{beamer}

\usepackage{../../style/lmu-lecture}
% Defines macros and environments
\input{../../style/common.tex}

\title{Applied Machine Learning}
% \author{LMU}
%\institute{\href{https://compstat-lmu.github.io/lecture_iml/}{compstat-lmu.github.io/lecture\_iml}}
\date{}

\begin{document}

\newcommand{\titlefigure}{figure/performance_vs_interpretability}
\newcommand{\learninggoals}{
\item What to watch out for in practice
\item Pros and cons of performance metrics
\item Advanced evaluation practices}

\lecturechapter{Applied Performance Evaluation}
\lecture{Applied Machine Learning}


\begin{frame}{General Things to Consider in Practice}

\begin{itemize}
    \item Stakeholders need to understand metrics
    \item Interpretation important? Easy to understand?
    \item What matters for the specific usecase, e.g., outlier sensitivity, cost sensitivity?
    \item Computational considerations?
    \item Type 3 error: Developing a model that answers the wrong question.
\end{itemize}
\end{frame}

\begin{frame}{What Affects Performance?}
% https://vuquangnguyen2016.files.wordpress.com/2018/03/applied-predictive-modeling-max-kuhn-kjell-johnson_1518.pdf

\begin{itemize}
    % \item Take this from Applied Predictive Modeling book
    \item Model type
    \item Non-informative features
    \item Pre-processing of features
    \item Measurement errors / noise in the outcome
    \item Measurement errors / noise in the features
    \item Sample size
    \item What kind of data? (algorithm-dependent)
\end{itemize}
\end{frame}


\begin{frame}{Calibration versus Discrimination}

\begin{itemize}
    \item Related to what we are trying to measure (type 3 error)
    \item Example with disease modeling
\end{itemize}

\end{frame}

\begin{frame}{Established Performance Metrics}

\begin{itemize}
    \item List and motivate established performance metrics here
    \item E.g., for \textbf{classification}: accuracy, true and false positive rate, precision and recall, F-score, AUC, Brier score
    \item E.g., for \textbf{regression}: MSE, RMSE, MAE, R2
    \item Describe typical use case and why metric works
    \item Example: Brier Score rewards model confidence (i.e. give a high probability to a true outcome) and punishes the reverse.
\end{itemize}
\end{frame}

\begin{frame}{Advanced Performance Metrics}

\begin{itemize}
    \item Dive into esoteric stuff here: Youden, Cohens-Kappa, Matthews correlation
    \item Motivate each advanced metric, do not just list formula
    \item Example: "MCC produces a high score only if the prediction obtained good results in all of the confusion matrix categories (TP, FP, TN, FN), proportionally both to the size of positive elements and the size of negative elements in the dataset"
\end{itemize}
\end{frame}

% \begin{frame}{Lift Curves}

% \begin{itemize}
%     \item Lift charts (Ling and Li 1998) are a visualization tool for assessing the ability of a model to detect events in a data set with two classes. 
%     \item They rank the samples by their scores and determine the cumulative event rate as more samples are evaluated. In the optimal case, the M highest-ranked samples would contain all M events.
% \end{itemize}
% \end{frame}

\begin{frame}{Lift Charts}

\begin{itemize}
\setlength\itemsep{1em}
    \item Lift charts visualize a model's ability to detect events.
    \item Samples are ranked according to their scores (probability for an event to take place). We then evaluate the cumulative event rate.
    \item In the optimal case, we would suspect the $m$ highest ranked samples to contain all $m$ events.
    \item For a non-informative model, the fraction $\frac{k}{n}$ of the highest ranked samples contains $k$ events on average.
    \item For a more informative model, the fraction $\frac{k}{n}$ of the highest ranked samples contains more than $k$ events.
\end{itemize}
\end{frame}



\begin{frame}{Lift Chart Example}

\begin{figure}
    \centering
    % \includegraphics{}
    \caption{Caption}
    \label{fig:my_label}
\end{figure}

\end{frame}


\begin{frame}{Pseudo Code to Create Lift Chart}

To construct a lift chart, we follow a general recipe:
\begin{itemize}
\setlength\itemsep{1em}
    \item Predict target values for a sample that was not used during training with known target values.
    \item Compute fraction of true events in the
entire data set.
    \item Rank data by scores (probability of event taking place).
    \item For every unique score,  
    For each unique class probability value, calculate the percent of true events
in all samples below the probability value.
\end{itemize}
\end{frame}

\begin{frame}{Calibration Curves}

\begin{itemize}
    \item Calibration curves (also known as reliability diagrams) compare how well the probabilistic predictions of a binary classifier are calibrated. It plots the true frequency of the positive label against its predicted probability, for binned predictions.
\end{itemize}
\end{frame}


\endlecture
\end{document}
